% Lösungen Größen und Einheiten (zusammenbau v0.4, Heft)
\einheitenkopf{Größen und Einheiten}

\erg{53}{a) Meter \quad b) $70$ dm \quad c) $3$ cm}
\erg{54}{a) $0{,}04$ m $= 4$ cm, also $0{,}04$ m $<$ $40$ cm \quad b) $2{,}7$ m $= 270$ cm, also $2{,}7$ m $>$ $27$ cm \quad c) $0{,}07$ kg $= 70$ g, also gilt $=$ \quad d) $55 \cdot 2{,}54$ cm $= 139{,}7$ cm $\approx 1{,}40$ m \quad e) $6 \cdot 1{,}61$ km $= 9{,}66$ km $\approx 9{,}7$ km}
\erg{55}{a) Zeitpunkt \quad b) $180$ s \quad c) $4$ min}
\erg{56}{a) $5{,}5 \cdot 60 = 330$ min \quad b) $6{,}5 \cdot 60 = 390$ min \quad c) $135$ min \quad d) 9:47 Uhr $+ 3$ h $=$ 12:47 Uhr, $+ 25$ min $=$ 13:12 Uhr \quad e) 16:52 Uhr $+ 1$ h $=$ 17:52 Uhr, $+ 43$ min $=$ 18:35 Uhr}
\erg{57}{a) hundert \quad b) $300$ dm² \quad c) $40\,000$ cm²}
\erg{58}{a) $1{,}25$ l $= 1\,250$ ml; $1\,250 : 200 = 6$ Rest $50$ – also $6$ volle Gläser \quad b) $4{,}2$ l $= 4\,200$ ml; $4\,200 : 750 = 5$ Rest $450$ – also $5$ volle Flaschen \quad c) $2{,}7$ l $= 2\,700$ ml; $2\,700 : 250 = 10$ Rest $200$ – also $10$ volle Kellen \quad d) $75$ m³ $= 75\,000$ l; $75\,000 : 12\,500 = 6$ h \quad e) $36$ m³ $= 36\,000$ l; $36\,000 : 800 = 45$ min}
\erg{59}{a) $1{,}2$ m \quad b) $1{,}4 + 0{,}65 = 2{,}05$ km \quad c) $8$ cm}
\erg{60}{a) $V = 3\,500 : 8{,}4 \approx 416{,}7$ cm³ \quad b) $V = 1\,800 : 11{,}3 \approx 159{,}3$ cm³ \quad c) $V = 620 : 7{,}3 \approx 84{,}9$ cm³ \quad d) Fläche $(2 \cdot 0{,}18 + 4 \cdot 0{,}12) \cdot 2 = 1{,}68$ m²; Bedarf $1{,}68 : 8 = 0{,}21$ l $= 210$ ml; $125$ ml reichen nicht, $250$ ml reichen – Max hat recht \quad e) Fläche $4 \cdot 0{,}9 \cdot 2 = 7{,}2$ m²; Bedarf $7{,}2 : 12 = 0{,}6$ l $= 600$ ml; $600$ ml $< 750$ ml – Ina hat recht \quad f) Fläche $(0{,}64 + 0{,}36) \cdot 2 \cdot 2 = 4$ m²; Bedarf $4 : 9 \approx 0{,}44$ l $\approx 444$ ml; $444$ ml $< 500$ ml – Jan hat recht \quad g) Fehler bei der Umrechnung: $1$ € $= 100$ ct, also $744$ € $= 74\,400$ ct, nicht $744\,000$ ct; $74\,400 : 120\,000 = 0{,}62$ ct – jeder Kilometer wird nur etwa $0{,}6$ ct teurer \quad h) Fehler bei der Umrechnung: $1$ € $= 100$ ct, also $396$ € $= 39\,600$ ct, nicht $396\,000$ ct; $39\,600 : 90\,000 = 0{,}44$ ct – pro Kilometer sind es nur etwa $0{,}4$ ct mehr}
